\documentclass[10pt,a4paper]{amsart}
\usepackage[margin=19mm]{geometry}
\usepackage{amsmath,amssymb,mathtools}
\usepackage[scr=rsfs,cal=euler]{mathalfa}
\usepackage{xfrac}
\usepackage[unicode,hidelinks,bookmarks=false]{hyperref}
\newcommand{\ie}{i.e.\ }
\newcommand{\cf}{cf.\ }
\newcommand{\resp}{respectively\ }
\newcommand{\Subsection}{\subsection}
\providecommand{\nobreakdash}{\nobreak\hbox{-}}
\setlength{\emergencystretch}{2em}
\multlinegap0pt
\usepackage{bm}
\usepackage{bbm}
\usepackage{dsfont}
\usepackage{xspace}
\usepackage{cancel}
\usepackage{mathtools}
\usepackage{cleveref}
\usepackage{amsthm}
\makeatletter
\@ifundefined{theorem}{\newtheorem{theorem}{Theorem}}{}
\@ifundefined{proposition}{\newtheorem{proposition}[theorem]{Proposition}}{}
\makeatother
\providecommand{\siff}{\Leftrightarrow}
\title[Eckstein-Ferris-Pennanen-Robinson duality revisited: paramon]{Eckstein-Ferris-Pennanen-Robinson duality revisited: paramonotonicity, total Fenchel-Rockafellar duality, and the Chambolle-Pock operator}
\author{Heinz H. Bauschke, Walaa M. Moursi, Shambhavi Singh}
\date{Source: arXiv:2412.11880v3 (CC BY 4.0)}
\begin{document}
\maketitle

\noindent\emph{Source and licence.} Eckstein-Ferris-Pennanen-Robinson duality revisited: paramonotonicity, total Fenchel-Rockafellar duality, and the Chambolle-Pock operator. Version arXiv:2412.11880v3 (CC BY 4.0), distributed under the Creative Commons Attribution 4.0 International licence, as recorded in the arXiv licence metadata for this exact version and shown on the abstract page https://arxiv.org/abs/2412.11880v3. Full rights notice: licenses/arxiv-2412.11880-CC-BY-4.0.txt.\par\smallskip

\noindent\textit{Editorial scope.} Counted unit: the Proposition whose source label is \texttt{p:241120b} in the article (source0.tex lines 596--623, source label \texttt{p:241120b}), delivered together with the article's own complete original proof of it (source0.tex lines 624--688, one \texttt{proof} environment of 65 lines). No mathematics is written, repaired or re-derived: statement and proof are the article's own text line for line. Required context retained uncounted: none. Every definition, display and label of those lines is retained unchanged. Omitted from this package and not redistributed: 1--595, 689--2592. Nothing inside a retained excerpt is omitted or altered: every retained body is delivered exactly as the article's lines are written, so no omission receipt of any kind accompanies this package. Citations: the retained excerpts carry no citation command at all, so no bibliography entry of the article is transcribed and no reference is redistributed. Reference adaptation: none was needed and none was made; every reference inside the delivered unit resolves inside it, so no unresolved reference or citation is produced. Printed slips kept literal: no printed word, symbol, hypothesis or number of the counted statement or of its proof was altered. Layout and class: the article's own class is \texttt{article} with packages including mathpazo, tikz, todonotes, soul, nameref, mathtools, empheq, comment, enumitem, hyperref, optional, xcolor, float, graphicx; the wrapper is the lane's pinned \texttt{amsart} editorial wrapper at 10pt on A4, which restates the theorem-like environments the retained text needs and adds the article's own macro definitions that the retained text needs (\texttt{\textbackslash siff}), with extra wrapper packages (bm, bbm, dsfont, xspace, cancel, mathtools, cleveref). The delivered numbering is the unit's own. Assets: no retained span contains a figure environment, a graphics-inclusion command, a PDF-page inclusion, a table or a TikZ picture; no image file is copied into this package, the package declares no asset dependency and no image file is redistributed with it. Licence: version arXiv:2412.11880v3 of this article is distributed under the Creative Commons Attribution 4.0 International licence, as recorded in the arXiv licence metadata for this exact version and shown on the abstract page https://arxiv.org/abs/2412.11880v3. A full rights notice is delivered at licenses/arxiv-2412.11880-CC-BY-4.0.txt. Characters: every retained excerpt is pure ASCII, so the delivered engine, XeLaTeX with its default Latin Modern text font, reports no missing character.
\begin{proposition}
\label{p:241120b}
Let $(x,y)\in X\times Y$. 
Then the following hold:
\begin{enumerate}
\item 
\label{p:241120b1}
$y\in K_x$ 
$\siff$ $x\in Z_y$. 
\item 
\label{p:241120b1.5}
$L^*(K_x) = (-Ax)\cap L^*BLx$ and 
$L(Z_y)= (B^{-1}y)\cap LA^{-1}(-L^*y)$.
\item 
\label{p:241120b2}
$K_x\neq\varnothing$
$\siff$ $x\in Z$, and
$Z_y\neq\varnothing$
$\siff$ $y\in K$.
\item 
\label{p:241120b3}
$K_x\subseteq K$ and 
$Z_y \subseteq Z$. 
\item 
\label{p:241120b4}
$K_x$ and $Z_y$ are closed and convex. 
\end{enumerate}
\end{proposition}

\begin{proof}
We have the equivalences
\begin{subequations}
\begin{align}
y\in K_x
&\siff
y\in -L^{-*}Ax \land y \in BLx\\
&\siff
-L^*y \in Ax \land Lx \in B^{-1}y\\
&\siff
x \in A^{-1}(-L^*y) \land x \in L^{-1}B^{-1}y\\
&\siff
x\in Z_y,
\end{align}
\end{subequations}
and this proves \cref{p:241120b1}.

We now prove \cref{p:241120b1.5}:
First, suppose that $v\in K_x$, i.e., 
$v\in-L^{-*}Ax$ and $v\in BLx$.
It follows that $L^*v\in -Ax$ and $L^*v\in L^*BLx$, i.e., 
$L^*v \in (-Ax)\cap L^*BLx$. 
We've verified that $L^*(K_x)\subseteq (-Ax)\cap L^*BLx$. 
Conversely, suppose that $u\in (-Ax)\cap L^*BLx$. 
Then there exists $v\in BLx$ such that $u=L^*v$. 
Hence $-L^*v\in Ax$ and so $v\in -L^{-*}x$. 
Thus $v\in K_x$ and so $u = L^*v \in L^*(K_x)$. 
Altogether, we have shown that $L^*(K_x) = (-Ax)\cap L^*BLx$. 
The proof of the identity 
$L(Z_y)= (B^{-1}y)\cap LA^{-1}(-L^*y)$ is similar. 


Turning to \cref{p:241120b2} and \cref{p:241120b3}, 
suppose that 
$K_x\neq\varnothing$, and let $k\in K_x$. 
It follows that 
\begin{equation}
k\in -L^{-*}Ax \land k \in BLx.
\end{equation}
First, this implies that 
$-L^*k \in Ax \land L^*k\in L^*BLx$, which upon adding, yields 
$x \in Z$. 
Second, this also implies 
$x \in A^{-1}(-L^*k) \land Lx \in B^{-1}k$. 
Hence $-Lx \in -LA^{-1}(-L^*k) \land Lx \in B^{-1}k$. 
Upon adding, we learn that $k\in K$. 
We have shown that $K_x \subseteq K$ and this is also trivially true 
if $K_x=\varnothing$. 

Now assume $x\in Z$. 
Then there exists $y\in BLx$ such that 
$-L^*y\in Ax$, i.e., $y\in -L^{-*}Ax$.
It follows that $y \in K_x$ and so $K_x\neq\varnothing$. 

The remaining statements in \cref{p:241120b2} and \cref{p:241120b3}
are proved similarly.

Because $A$ and $B$ are maximally monotone, 
the sets 
$Ax$ and $B(Lx)$ are convex and closed. 
The linearity and continuity of $-L^*$ then 
yields the convexity and closedness of $-L^{-*}Ax$.
Upon intersecting, we obtain  convexity and closedness $K_x$. 
The set $Z_k$ is treated similarly. 
\end{proof}

\end{document}
